% ============================================================
%  Abstract
% ============================================================
\chapter*{Abstract}
\addcontentsline{toc}{chapter}{Abstract}
\thispagestyle{plain}

Groundwater is the primary source of freshwater for irrigation and domestic
use across the semi-arid Pothohar Plateau of northern Punjab, Pakistan, yet
its spatial distribution remains poorly mapped at a regional scale. This
thesis develops a machine learning (ML) framework for groundwater potential
mapping (GWPM) that extends the knowledge-based weighted-overlay approach of
Waqas et al.\ \cite{waqas2022chakwal} from a single district to the full
five-district plateau (Chakwal, Rawalpindi, Attock, Jhelum, and Mianwali;
$\approx$\SI{25000}{\kilo\meter\squared}).

Nine geo-environmental features --- elevation, slope, aspect, topographic
wetness index (TWI), NDVI, NDWI, long-term rainfall, soil texture, and land
surface temperature --- were derived from Landsat~8, SRTM, CHIRPS, MODIS, and
SoilGrids using Google Earth Engine at 30~m resolution, yielding \npixels{}
valid pixels. A knowledge-based weighted overlay generated three-class
groundwater potential labels (Low, Medium, High), from which a balanced
training set of 450{,}000 samples was drawn. Four classifiers --- Random
Forest, XGBoost, LightGBM, and a soft-voting ensemble --- were trained and
evaluated.

XGBoost achieved the best performance (test accuracy 99.12\%, macro
$F_1=0.991$, Cohen's $\kappa=0.987$, ROC-AUC $=0.9998$) and was applied to
the full plateau, classifying 20\% of the area as High, 40\% as Medium, and
40\% as Low potential. Because the training labels and model inputs share
the same feature space, internal accuracy primarily measures reproducibility
of the overlay; the model's real-world credibility was therefore established
through \emph{spatial-realism validation} against independent evidence:
(i)~4/4 published alluvial groundwater sites were correctly classified as
Medium or High; (ii)~alluvial rain-gauge stations showed significantly higher
predicted potential than non-alluvial stations (Mann--Whitney $p=0.023$);
(iii)~district-level predictions followed the expected hydrogeological
gradient (Spearman $\rho=1.00$, Rawalpindi 61\% High vs.\ Chakwal 4\% High);
and (iv)~the High-potential fraction (20\%) was consistent with two
independent 2025 AHP studies of the same region.

The resulting 30~m groundwater potential map identifies the Jhelum, Soan, and
Indus alluvial corridors as priority zones for well development. The study
demonstrates that ML can replace manual weighted overlay while retaining
hydrogeological interpretability, and provides a reproducible, scalable
framework for regional groundwater assessment in data-scarce settings.

\vspace{0.4cm}
\noindent\textbf{Keywords:} Groundwater potential mapping; Machine learning;
XGBoost; Random Forest; Google Earth Engine; Remote sensing; Pothohar Plateau;
Pakistan; Spatial validation.

\cleardoublepage
